\section{先備知識(波粒二像性與光譜)}
\subsection{波粒二像性與光能}

光同時擁有波動性與粒子性，兩者分別具備以下性質:
\begin{itemize}
    \item \textbf{波動性}
    \begin{enumerate}
        \item 光為電磁波，在真空中的速度為 $3\times10^8\ \mathrm{m/s}$。
        \item 光會產生干涉與繞射等波動現象。
        \item 光具有波長與頻率，兩者關係為$c=\lambda \nu$\\
        (化學上頻率常以$\nu$表示，物理則是f，但其實是一樣的東西。)
    \end{enumerate}

    \item \textbf{粒子性}
    \begin{enumerate}
        \item 光由光子（photon）組成。
        \item 每個光子的能量為$E=h\nu$。
        \item 光的粒子性可由光電效應等實驗證明(參見選修物理第五冊)。
    \end{enumerate}
\end{itemize}

\subsection{光譜}

\begin{trick}
定義:\\
各種不同頻率的光在在感光元件上形成的亮線或暗線分佈。
\end{trick}
光譜分為三種:
\begin{itemize}
    \item \textbf{連續光譜}：包含所有波長，沒有間斷。
    \item \textbf{發射光譜}：原子激發後釋放能量形成的亮線光譜。
    \item \textbf{吸收光譜}：連續光通過氣體時，特定波長被吸收而形成暗線。
\end{itemize}
而在本小節主要探討的是氫原子的\underline{發射光譜}。




%$\frac{1}{\lambda}=R_H(\frac{1}{n_L^2}-\frac{1}{n_H^2})$